\documentclass[10pt,a4paper]{amsart}
\usepackage[margin=19mm]{geometry}
\usepackage{amsmath,amssymb,mathtools}
\usepackage[scr=rsfs,cal=euler]{mathalfa}
\usepackage{xfrac}
\usepackage[unicode,hidelinks,bookmarks=false]{hyperref}
\newcommand{\ie}{i.e.\ }
\newcommand{\cf}{cf.\ }
\newcommand{\resp}{respectively\ }
\newcommand{\Subsection}{\subsection}
\providecommand{\nobreakdash}{\nobreak\hbox{-}}
\setlength{\emergencystretch}{2em}
\multlinegap0pt
\usepackage{bm}
\usepackage{bbm}
\usepackage{dsfont}
\usepackage{xspace}
\usepackage{cancel}
\usepackage{mathtools}
\usepackage{amsthm}
\makeatletter
\@ifundefined{lem}{\newtheorem{lem}{Lemma}}{}
\makeatother
\newcommand{\C}{\mathbb{C}}
\newcommand{\Es}{\mathscr{E}}
\newcommand{\Fb}{\mathbf{F}}
\newcommand{\Gb}{\mathbf{G}}
\newcommand{\Gr}{\mathsf{Gr}}
\newcommand{\Ls}{\mathscr{L}}
\newcommand{\Mean}{\mathscr{M}}
\newcommand{\N}{\mathbb{N}}
\newcommand{\R}{\mathbb{R}}
\newcommand{\T}{\mathbb{T}}
\newcommand{\Ub}{\mathbf{U}}
\newcommand{\Us}{\mathscr{U}}
\newcommand{\be}{\begin{equation}}
\newcommand{\ee}{\end{equation}}
\newcommand{\eu}{\mathrm{e}}
\newcommand{\id}{I}
\newcommand{\weakto}{\rightharpoonup}
\renewcommand{\geq}{\geqslant}
\renewcommand{\ker}{\mathrm{Ker}}
\title[The half-wave maps equation on $\mathbb{T}$: Global well-pos]{The half-wave maps equation on $\mathbb{T}$: Global well-posedness in $H^{1/2}$ and almost periodicity}
\author{Patrick G\'erard, Enno Lenzmann}
\date{Source: arXiv:2602.20895v3 (CC BY 4.0)}
\begin{document}
\maketitle

\noindent\emph{Source and licence.} The half-wave maps equation on $\mathbb{T}$: Global well-posedness in $H^{1/2}$ and almost periodicity. Version arXiv:2602.20895v3 (CC BY 4.0), distributed under the Creative Commons Attribution 4.0 International licence, as recorded in the arXiv licence metadata for this exact version and shown on the abstract page https://arxiv.org/abs/2602.20895v3. Full rights notice: licenses/arxiv-2602.20895-CC-BY-4.0.txt.\par\smallskip

\noindent\textit{Editorial scope.} Counted unit: the Lem whose source label is \texttt{lem:eigenfunctions} in the article (source0.tex lines 1791--1796, source label \texttt{lem:eigenfunctions}), delivered together with the article's own complete original proof of it (source0.tex lines 1798--1847, one \texttt{proof} environment of 50 lines). No mathematics is written, repaired or re-derived: statement and proof are the article's own text line for line. Required context retained uncounted: lem:commutator (lines 798--809); lem:UsSigma (lines 1485--1498); lem:GWP\_minor (lines 1697--1703). Every definition, display and label of those lines is retained unchanged. Omitted from this package and not redistributed: 1--797, 810--1484, 1499--1696, 1704--1790, 1797--1797, 1848--2309. Nothing inside a retained excerpt is omitted or altered: every retained body is delivered exactly as the article's lines are written, so no omission receipt of any kind accompanies this package. Citations: the retained excerpts carry no citation command at all, so no bibliography entry of the article is transcribed and no reference is redistributed. Reference adaptation: none was needed and none was made; every reference inside the delivered unit resolves inside it, so no unresolved reference or citation is produced. Printed slips kept literal: no printed word, symbol, hypothesis or number of the counted statement or of its proof was altered. Layout and class: the article's own class is \texttt{amsart} with packages including amsmath, amssymb, amsthm, mathtools, dsfont, bm, bbm, mathrsfs, color; the wrapper is the lane's pinned \texttt{amsart} editorial wrapper at 10pt on A4, which restates the theorem-like environments the retained text needs and adds the article's own macro definitions that the retained text needs (\texttt{\textbackslash C}, \texttt{\textbackslash Es}, \texttt{\textbackslash Fb}, \texttt{\textbackslash Gb}, \texttt{\textbackslash Gr}, \texttt{\textbackslash Ls}, \texttt{\textbackslash Mean}, \texttt{\textbackslash N}, \texttt{\textbackslash R}, \texttt{\textbackslash T}, \texttt{\textbackslash Ub}, \texttt{\textbackslash Us}, \texttt{\textbackslash be}, \texttt{\textbackslash ee}, \texttt{\textbackslash eu}, \texttt{\textbackslash geq}, \texttt{\textbackslash id}, \texttt{\textbackslash ker}, \texttt{\textbackslash weakto}), with extra wrapper packages (bm, bbm, dsfont, xspace, cancel, mathtools). The delivered numbering is the unit's own. Assets: no retained span contains a figure environment, a graphics-inclusion command, a PDF-page inclusion, a table or a TikZ picture; no image file is copied into this package, the package declares no asset dependency and no image file is redistributed with it. Licence: version arXiv:2602.20895v3 of this article is distributed under the Creative Commons Attribution 4.0 International licence, as recorded in the arXiv licence metadata for this exact version and shown on the abstract page https://arxiv.org/abs/2602.20895v3. A full rights notice is delivered at licenses/arxiv-2602.20895-CC-BY-4.0.txt. Characters: every retained excerpt is pure ASCII, so the delivered engine, XeLaTeX with its default Latin Modern text font, reports no missing character.
\begin{lem} \label{lem:commutator}
For $\Ub \in L^\infty(\T; \Ls(E))$ and $\Fb \in L^2_+(\T;E)$, it holds
$$
[S^*, T_\Ub] \Fb = S^* T_{\Ub}(\Mean(\Fb)).
$$

In particular, if  $E=\C^{d \times d}$ and $\Ub \in L^\infty(\T; \C^{d \times d})$, the identity above reads
$$
[S^*, T_\Ub] \Fb = (S^* \Pi \Ub)\Mean(\Fb),
$$
with the backward shift $S^*$ acting on $\Pi \Ub \in L^2_+(\T, \C^{d \times d})$.
\end{lem}

\begin{lem} \label{lem:UsSigma}
The linear map $\Us_\Sigma : H \to L^2_+(\T;E)$ is a contraction. For all $\Fb \in H$, the following properties hold.
\begin{enumerate}
\item[(i)] We have the identity
$$
\| \Us_\Sigma \Fb \|_{L^2_+}^2 = \| \Fb \|^2 - \lim_{N \to \infty} \| (\Sigma^*)^N \Fb \|^2.
$$
\item[(ii)] $\Us_\Sigma(\Fb) = \Fb$ whenever $\Fb \in E$.
\item[(iii)] The following intertwining relations hold
$$
S \Us_\Sigma = \Us_\Sigma \Sigma \quad \mbox{and} \quad S^* \Us_\Sigma = \Us_{\Sigma} \Sigma^* \quad \mbox{on $H$}.
$$ 
\end{enumerate}
\end{lem}

\begin{lem} \label{lem:GWP_minor}
Let $\Ub_0 \in H^{1/2}(\T; \Gr_k(\C^d))$. Suppose that the sequence $\Ub_{0,n} \in \mathcal{R}at(\T; \Gr_k(\C^d))$ with $n \in \N$ satisfies $\Ub_{0,n} \to \Ub_0$ in $H^{1/2}$ and let $t_n \to t$. Then 
$$
\mbox{$\Us_n(t_n) \Pi \Ub_{0,n} \to \Us(t) \Pi \Ub_0$ in $L^2$} \quad \mbox{and} \quad \mbox{$\Us_n(t_n) \Pi \Ub_{0,n} \weakto \Us(t) \Pi \Ub_0$ in $H^{1/2}$}
$$
In particular, the limit $\Us(t) \Pi \Ub_0$ is independent of the chosen sequence of rational initial data $\Ub_{0,n} \in \mathcal{R}at(\T; \Gr_k(\C^d))$. 
\end{lem}

\begin{lem} \label{lem:eigenfunctions}
Let $t \in \R$ be given. Suppose $\Fb \in L^2_+ \setminus \{ 0 \}$ is an eigenfunction for $T_{\Ub_0}$ with eigenvalue $\mu \in \sigma(T_{\Ub_0})$. Then $\Fb(t) := \Us(t) \Fb$ satisfies
$$
T_{\Ub(t)} \Fb(t) = \mu \Fb(t) \quad \mbox{and} \quad \Fb(t) \neq 0.
$$
\end{lem}

\begin{proof}
We divide the proof into the following steps.

\medskip
\textbf{Step 1.} As before, let $\Ub_{0,n} \in \mathcal{R}at(\T; \Gr_k(\C^d))$ be a sequence of rational data with $\Ub_{0,n} \to \Ub_0$ in $H^{1/2}$. Likewise, we denote by $\Ub_n(t)$ the corresponding solutions $C(\R; H^\infty)$ of (HWM) with $\Ub_n(0) = \Ub_{0,n}$ and we let $\Us_n(t) : L^2_+ \to L^2_+$ denote the corresponding sequence of unitary maps. From the Lax evolution for rational solutions we infer that
\be \label{eq:Lax_eigen}
T_{\Ub_n(t)} \Us_n(t) \Fb = \Us_n(t) T_{\Ub_0} \Fb = \mu \Us_n(t) \Fb
\ee
where $\Fb \in L^2_+ \setminus \{ 0 \}$ solves $T_{\Ub_0} \Fb = \mu \Fb$. Note that $\Us_n(t) \Fb \weakto \Us(t) \Fb$ in $L^2$ analogous to the proof of Lemma \ref{lem:GWP_minor}. Moreover, from $\Ub_n(t) \to \Ub(t)$ in $L^2$ we easily deduce that $T_{\Ub_n(t)} \to T_{\Ub(t)}$ strongly as operators. By this fact and the self-adjointness of $T_{\Ub_n(t)}$ and $T_{\Ub(t)}$, it is elementary to check that $T_{\Ub_n(t)} \Us_n(t) \Fb \weakto T_{\Ub(t)} \Us(t) \Fb$. By passing to the limit $n \to \infty$ in \eqref{eq:Lax_eigen}, we obtain that $\Fb(t) := \Us(t) \Fb$ satisfies
$$
T_{\Ub(t)} \Fb(t) = \mu \Fb(t) .
$$

\medskip
\textbf{Step 2.} We now claim that 
$$
\Fb(t) \neq 0 .
$$
Recalling Lemma \ref{lem:UsSigma}, we see that $\Fb(t) = \Us(t) \Fb \neq 0$ provided that $\Mean(\Fb) \neq 0$. Hence it remains to discuss the case when $\Mean(\Fb) = 0$ for the rest of the proof.

Suppose now that $\Fb \neq 0$ with $\Mean(\Fb)=\widehat{\Fb}_0=0$. Since $\Fb \neq 0$, we can define 
$$
n := \min \{ k \geq 1 : \Mean((S^*)^k \Fb) = \widehat{\Fb}_k \neq 0 \}.
$$
Next, we consider the eigenspace
$$
\Es_{\mu} = \ker ( T_{\Ub_0} - \mu \id).
$$
As an eigenspace for the Toeplitz operator $T_{\Ub_0}$, we recall that $\Es_\mu$ must be {\em nearly $S^*$-invariant}, i.e., 
$$
\Gb \in \Es_\mu \quad \mbox{and} \quad \Mean(\Gb) = 0 \quad \Rightarrow \quad S^* \Gb \in \Es_\mu,
$$
which is a direct consequence of the commutator identity in Lemma \ref{lem:commutator}. Therefore, 
$$
(S^*)^j \Fb \in \Es_{\mu} \quad \mbox{for $j=0, \ldots, n$}.
$$ 
Now we recall the intertwining identity from Lemma \ref{lem:UsSigma}, which gives us
\be \label{eq:intertwine}
S^* \Us(t) = \Us(t) \eu^{-it T_{\Ub_0}} S^*.
\ee
By iterating this identity $n$ times and applying it to the vector $\Fb$, we obtain
$$
(S^*)^n \Us(t) \Fb = \Us(t) (\eu^{-it T_{\Ub_0}} S^*)^n \Fb.
$$
Since $\eu^{-it T_{\Ub_0}} \Gb = \eu^{-it \mu} \Gb$ for $\Gb \in \Es_\mu$, we find $(\eu^{-it T_{\Ub_0}} S^*)^n \Fb= \eu^{-in \mu t} (S^*)^n \Fb$, whence it follows that
$$
(S^*)^n \Us(t) \Fb= e^{-in \mu t} \Us(t) (S^*)^n \Fb \neq 0,
$$
using that $\Mean((S^*)^n \Fb) \neq 0$ and the fact that $\Us(t) \Gb \neq 0$ if $\Mean(\Gb) \neq 0$ by Lemma \ref{lem:UsSigma}. Thus we deduce that $\Fb(t) = \Us(t) \Fb \neq 0$ must hold.
\end{proof}

\end{document}
